\documentclass[11pt]{article}
\title{How the Codynamic Book Machine Works}
\author{Arthur Petron}
\usepackage[margin=1in]{geometry}
\usepackage{graphicx}
\usepackage{amsmath,amssymb,mathtools}
\usepackage{tikz}
\usetikzlibrary{positioning,arrows.meta}
\usepackage{hyperref}
\graphicspath{{media/}{media/diagrams/}{images/}{artwork/}}

\begin{document}

\section*{Conversational Intake}

Conversational intake is the system's first structured conversation with the author. Its purpose is to turn a vague project idea into a book-shaped brief that can be normalized, validated, and handed to the rest of the machine. Rather than asking the author to produce a complete outline up front, the intake process asks a small set of targeted questions that capture the minimum information needed to begin: the title, the purpose of the work, the intended audience, the reader takeaway, the desired voice, the genre, the design specification, and the expected artifact structure.

These questions are not arbitrary metadata. Each one constrains a different part of the downstream workflow. The title gives the project a stable identity. The purpose states what the book is trying to accomplish. The audience and reader takeaway define who the book is for and what success looks like from the reader's point of view. Voice and genre establish the rhetorical and stylistic frame so that later drafting can remain consistent. The design specification describes how the book should be shaped as an artifact, including any structural or presentation requirements that affect generation and compilation. Artifact structure tells the system what kinds of outputs are expected and how they should be organized.

Taken together, the question bank functions as a compact schema for author intent. It is not trying to capture every detail of the finished book; it is trying to capture the minimum set of commitments that make the rest of the pipeline reliable. A title without purpose is only a label. A purpose without audience is only an aspiration. An audience without reader takeaway does not yet define success. The intake conversation therefore asks for these fields in combination so that the system can build a coherent project model rather than a loose collection of preferences.

The question bank also separates discovery from drafting without separating them completely. The author is not forced to solve the whole book in one pass, but the system still receives enough information to normalize the project into machine-readable fields, check for missing pieces, and seed the outline and section workflow that follows. In practice, conversational intake behaves like a guided interview and a schema-filling step at the same time. It is conversational because it helps the author think through the project. It is operational because each answer becomes part of the canonical work object.

That dual role matters for the rest of the machine. Once the intake answers are recorded, later agents can rely on them as stable context instead of repeatedly inferring intent from free-form prose. The outline can be shaped around the stated purpose. Drafting can follow the declared voice and genre. Design and assembly can respect the artifact structure from the start. In this way, the intake stage is not merely administrative overhead; it is the point where author intent becomes operational structure.

The reader takeaway from this stage is simple: before the machine can write, it must know what it is writing, for whom, and in what form. Conversational intake supplies that information in a controlled way, so the rest of the pipeline can operate on a coherent project rather than an ambiguous prompt.


\end{document}
